\section{Ctl-1: one-shot collateral geometry}

Ctl-1 is the smallest test in the paper. The activation state is
\[
h=[x_1,x_2,x_1x_2].
\]
The target readout is
\[
z=\beta_1x_1+\beta_2x_2+\gamma x_1x_2,
\]
with default $(\beta_1,\beta_2,\gamma)=(0.35,0.25,1.15)$. The nuisance
readout is
\[
n(w)=(1-w)[0.90,-0.25,0]+w[0,0,1].
\]
The first-order estimator uses only $x_1$ and $x_2$. The lifted estimator
also uses $x_1x_2$. If $\widetilde d$ is a raw direction, the original
fixture scales it so that
\[
d=\frac{g_0}{\nabla z(h_0)^\top\widetilde d}\widetilde d,
\qquad
\nabla z(h_0)^\top d=g_0.
\]
The one-shot controller applies
\[
h^+=h_0+ud,
\qquad
u=K(y^\star-\widehat z_0),
\]
and collateral is
\[
C=\left|n(w)^\top(h^+-h_0)\right|.
\]

This construction is a coordinate-level intervention. After the update, its
third coordinate need not equal the product of its first two coordinates.
Calling the third coordinate an ``interaction'' does not make the post-edit
state nonlinear. The primary result is therefore an omitted-coordinate
geometry result: a readout that cannot use one coordinate directs the edit
through the coordinates it can see.

\subsection{Null and interaction sweeps}

Setting $\gamma=0$ removes the omitted interaction coordinate from the target.
Across ten seeds, first-order and lifted arms then agree: target MSE is $0.0690$
for both, and collateral is $0.7318$ and $0.7317$, respectively. The task
does not reward the lifted estimator when the target is first order.

With the nuisance in the main-effect subspace, increasing $\gamma$ makes the
first-order arm more collateral-heavy. At $\gamma=0.5$, it obtains $98.0\%$
of the lifted arm's target improvement while paying about $2.1\times$ the
collateral. At larger interaction strengths, the target errors also separate,
so the moderate point states the collateral result more cleanly.

\begin{figure}[t]
  \centering
  \includegraphics[width=0.92\linewidth]{figure_01_ctl1_analytic_interaction_sweep.png}
  \caption{Analytic Ctl-1 interaction sweep. The useful regime is where the
  first-order arm remains target-competitive but already pays more collateral.}
  \label{fig:ctl1-interaction}
\end{figure}

Moving the nuisance readout changes which direction pays the cost. Ratios above
one mean the first-order direction moves more collateral; ratios below one mean
the lifted direction does. The crossover in Figure~\ref{fig:ctl1-nuisance}
shows that collateral follows direction--readout overlap, not the observer name.

\begin{figure}[t]
  \centering
  \includegraphics[width=0.92\linewidth]{figure_02_ctl1_analytic_nuisance_sweep.png}
  \caption{Analytic nuisance placement. The collateral burden crosses over as
  the nuisance moves from the main-effect coordinates toward the omitted
  coordinate.}
  \label{fig:ctl1-nuisance}
\end{figure}

\subsection{Base-coordinate robustness check}

We also edit the base coordinates and recompute the product exactly. For a
two-dimensional direction $q$,
\[
(x_1^+,x_2^+)=(x_1,x_2)+uq,
\qquad
h^+=[x_1^+,x_2^+,x_1^+x_2^+].
\]
At seed 0, the baseline target MSE is $0.718$. The lifted arm reduces it to
$0.490$, while the first-order arm increases it to $1.259$. The result does
not depend on freeing the third coordinate.

This check is still continuous. About $88\%$ of edited examples leave the
binary unit square in both arms. It therefore tests local mechanism geometry,
not a realizable token intervention.
